\documentclass[11pt]{article}

\usepackage[margin=1in]{geometry}
\usepackage{hyperref}
\usepackage{booktabs}
\usepackage{longtable}
\usepackage{array}
\usepackage{enumitem}
\usepackage{listings}
\usepackage{xcolor}
\usepackage{parskip}
\usepackage{titlesec}
\usepackage{seqsplit}

\hypersetup{colorlinks=true, linkcolor=blue, urlcolor=blue, citecolor=blue}

\lstset{
  basicstyle=\ttfamily\small,
  breaklines=true,
  columns=fullflexible,
  frame=single,
  backgroundcolor=\color{gray!8},
  showstringspaces=false
}

\newcommand{\code}[1]{\texttt{\seqsplit{#1}}}
\newcommand{\bslash}{\texttt{\seqsplit{\char`\\}}}

\titleformat{\section}{\large\bfseries}{\thesection}{0.6em}{}
\titleformat{\subsection}{\normalsize\bfseries}{\thesubsection}{0.6em}{}

\title{\textbf{Claude-powered paper audit pipeline --- design report}}
\author{}
\date{}

\begin{document}
\maketitle
\vspace{-2em}

\noindent\textbf{Repo:} \code{sednabcn/LLM-HypatiX-DEV} \\
\textbf{Scope:} audits \code{**/*\_patched.tex} against \code{hypatiax/data/results/}, using the Claude API, independent of the existing \code{ci\_paper\_audit.yml} gate.

\vspace{1em}
\hrule
\vspace{1em}

\section{Why this exists}

\code{ci\_paper\_audit.yml} already does deterministic, rule-based auditing of the paper against results data. This pipeline adds a second, LLM-driven layer on top, for two jobs a deterministic script can't do:

\begin{enumerate}
  \item \textbf{Numeric consolidation} --- when a paper claim's underlying results have multiple files (reruns, seeds, shards), decide the single authoritative value and \emph{why}, then propose (and in narrow cases, auto-commit) the fix.
  \item \textbf{The narrative ``open/flagging points'' report} --- a human-readable account of what's stale, what's missing, and what's actively blocked, on top of the raw findings JSON.
\end{enumerate}

It deliberately does \textbf{not} try to replicate the deep, cascading, multi-file investigations run manually (the \code{68\_74-issue.txt} / \code{leak\_report.tex} style sessions). Those stay human+Claude, outside CI. This pipeline's only relationship to that work is: \textbf{it reads the block list those sessions produce, and it never overwrites a number in a table/figure that list says is currently invalid.}

\section{Two-tier design}

\begin{longtable}{@{}p{2.8cm}p{6cm}p{6cm}@{}}
\toprule
 & \textbf{Tier 1 --- automated, in CI} & \textbf{Tier 2 --- human+Claude, outside CI} \\
\midrule
\endhead
What it is &
  \code{llm\_audit\_inspector.py}, one Claude call per claim, bounded &
  The kind of session in \code{68\_74-issue.txt} / \code{notes-nguyen12.txt} --- iterative, requests more files, hypotheses get overturned, scope legitimately expands \\
Trigger &
  push to \code{*\_patched.tex} / \code{hypatiax/data/results/**} / \code{audit/observations/**}, or manual dispatch &
  you, manually, whenever \\
Output &
  \code{logs/llm\_audit\_findings.json}, updates to \code{audit\_registry.json}, three forms of report &
  Incident docs like \code{leak\_report.tex}, filed back in via \code{audit/observations/*.json} (\code{type: "block"}) \\
Can it auto-commit? &
  Yes, but \textbf{only} a single-digit numeric substitution, mechanically re-verified (no LLM trust) &
  Never automatically --- always a human-reviewed patch \\
Relationship &
  \textbf{Reads} Tier 2's block list before doing anything; refuses to touch any claim inside a blocked table/figure &
  Produces the block list; the only way to remove a block is editing \code{audit\_registry.json} directly \\
\bottomrule
\end{longtable}

\section{Data flow}

\begin{lstlisting}
push to *_patched.tex / hypatiax/data/results/** / audit/observations/**
                |
                v
   +--------------------------+
   |  1. INSPECT               |  llm_audit_inspector.py
   |  - extract_claims()       |  reads: *_patched.tex,
   |  - merge audit_registry   |         hypatiax/data/results/**,
   |    (seed/tier2 preserved) |         audit/observations/*,
   |  - blocklist check        |         audit_registry.json (existing)
   |  - consolidate() per      |  calls: consolidate_results.py (Claude API)
   |    unblocked claim        |  writes: logs/llm_audit_findings.json,
   +--------------------------+          audit_registry.json (merged)
                |
                v
   +--------------------------+
   |  2. PATCH                 |  llm_patch_apply.py
   |  generate: findings ->    |  reads: logs/llm_audit_findings.json
   |    patch_manifest.yaml    |  writes: patch_manifest.yaml (proposed
   |    entries                |          entries, all kinds)
   |  apply-numeric: verify +  |  reads/writes: patch_manifest.yaml,
   |    apply ONLY             |          *_patched.tex (numeric_* only)
   |    numeric_consolidated/  |  auto-commits ONLY these, mechanically
   |    numeric_from_upload    |  re-verified line-by-line
   +--------------------------+
                |
                v
   +--------------------------+
   |  3. REPORT                |  generate_llm_audit_report.py
   |  writes all 3 forms       |  reads: logs/llm_audit_findings.json,
   +--------------------------+          audit_registry.json
                |
        +-------+----------------+
        v       v                v
 audit/llm_    logs/llm_audit_   <dir>/<file>_llm_audit.md
 audit_report  report_<run_id>.  (per *_patched.tex,
 .md (cumula-  md (artifact       committed next to it)
 tive,          only, not
 committed)     committed)
\end{lstlisting}

Separately, whenever a Tier-2 investigation produces a block:

\begin{lstlisting}
you file audit/observations/<name>.json  (type: "block")
                |
                v
next Tier-1 run's INSPECT step merges it into audit_registry.json
as a tier2_upload entry -> its `blocks` list is now enforced
\end{lstlisting}

\section{File inventory --- what goes where in the repo}

All paths below are relative to the repo root.

\begin{longtable}{@{}p{0.4cm}p{4.2cm}p{4.6cm}p{5.5cm}@{}}
\toprule
\# & \textbf{File} & \textbf{Repo destination} & \textbf{Purpose} \\
\midrule
\endhead
1 & \code{ci\_llm\_paper\_audit.yml} & \code{.github/workflows/} & The workflow itself: 3 jobs (\code{inspect} $\to$ \code{patch} $\to$ \code{report}), independent of \code{ci\_paper\_audit.yml} \\
2 & \code{llm\_audit\_inspector.py} & \code{scripts/patches/} & Tier-1 entry point. Extracts numeric claims from \code{*\_patched.tex}, tracks table/figure labels, checks the blocklist, calls \code{consolidate\_results.py} for unblocked claims, merges results into \code{audit\_registry.json} \\
3 & \code{consolidate\_results.py} & \code{scripts/patches/} & The \textbf{only} file that calls the Claude API. One bounded call per claim: given a section's result summary files, returns the authoritative value + method + confidence + provenance as strict JSON \\
4 & \code{llm\_patch\_apply.py} & \code{scripts/patches/} & Deterministic, no Claude calls. \code{generate} writes findings into \code{patch\_manifest.yaml} (existing schema); \code{apply-numeric} mechanically verifies and applies \emph{only} numeric single-token substitutions \\
5 & \code{generate\_llm\_audit\_report.py} & \code{scripts/patches/} & Writes the open/flagging-points report in all three requested forms (cumulative, timestamped, per-tex-file) \\
6 & \code{audit\_registry\_lib.py} & \code{scripts/patches/} & Shared library: load/save \code{audit\_registry.json}, compute the blocklist, merge Tier-1 findings without ever touching seed/\code{tier2\_upload} entries \\
7 & \code{seed\_audit\_registry.py} & \code{scripts/patches/} & One-time (idempotent) importer --- merges a seed file into \code{audit\_registry.json} \\
8 & \code{plan\_action\_seed.json} & \code{audit/seed/} & The actual 13 \code{PLAN-ACTION.txt} items + item 14 (the ground-truth leak from \code{leak\_report.tex}), pre-converted to the registry schema \\
9 & \code{README.md} (observations) & \code{audit/observations/} & Documents both upload paths (committed \code{.md}/\code{.json}, and \code{workflow\_dispatch} ad-hoc) and the \code{type: "block"} schema for filing Tier-2 findings \\
\bottomrule
\end{longtable}

\subsection*{Files this pipeline generates at runtime (not delivered --- created by the workflow)}

\begin{longtable}{@{}p{4.5cm}p{4.0cm}p{2.2cm}p{3.9cm}@{}}
\toprule
\textbf{File} & \textbf{Path} & \textbf{Committed?} & \textbf{Notes} \\
\midrule
\endhead
\code{audit\_registry.json} & repo root & Yes & merged and committed each run \\
\code{patch\_manifest.yaml} & repo root & Yes & existing file; this pipeline only appends \\
\code{logs/llm\_audit\_findings.json} & \code{logs/} & No & artifact only \\
\code{audit/llm\_audit\_report.md} & \code{audit/} & Yes & cumulative, overwritten each run \\
\code{logs/llm\_audit\_report\_<run\_id>.md} & \code{logs/} & No & artifact only \\
\code{<dir>/<texfile>\_llm\_audit.md} & next to each \code{*\_patched.tex} & Yes & one per patched tex file \\
\code{audit/incident\_reports/*.tex} & \code{audit/incident\_reports/} & manual & you commit these; pipeline only reads/links \\
\bottomrule
\end{longtable}

\section{Setup steps (in order)}

\begin{enumerate}
  \item Copy files 1--9 from Section~4 into the listed paths.
  \item \code{pip install anthropic pyyaml} (already in the workflow's install step --- nothing to do if only running in CI).
  \item Confirm \code{secrets.ANTHROPIC\_API\_KEY} is set at the repo level (it already is, per \code{ci\_paper\_audit.yml}).
  \item \textbf{Seed the registry once:}
\begin{lstlisting}
python3 scripts/patches/seed_audit_registry.py \
  --seed audit/seed/plan_action_seed.json \
  --registry audit_registry.json
git add audit_registry.json && git commit -m \
  "chore: seed audit_registry.json from PLAN-ACTION.txt"
\end{lstlisting}
  This immediately puts Tables 9--13, 15, \code{tab:hybrid\_all}, \code{tab:runtime}, \code{tab:timing}, and Figures 9--13 on the blocklist, and items \#4, \#10, \#12, \#13 in as \code{resolved} so Tier 1 doesn't re-flag them.
  \item If \code{audit/incident\_reports/leak\_report.tex} already exists, commit it now too --- the seed entry for item 14 already points \code{linked\_report} at that exact path.
  \item \textbf{Fix the one assumption made blind:} \code{resolve\_section\_to\_results()} in \code{llm\_audit\_inspector.py} assumes a \code{paper\_targets.json} schema of the form \code{\{"sections": [\{"section": ..., "experiment\_ids": [...]\}]\}}, which was invented for lack of the real file. Supply the real schema (or the equivalent lookup already inside \code{hypatia\_inspector.py}) to rewire that one function --- everything else is independent of it.
\end{enumerate}

\section{Known limitations (found via testing, not theoretical)}

\begin{itemize}[leftmargin=1.2em]
  \item \textbf{Label detection is regex-based, not a real tex parser.} \code{LABEL\_RE} catches \code{\textbackslash label\{tab:...\}}/\code{\textbackslash label\{fig:...\}}; \code{PROSE\_LABEL\_RE} catches literal ``Table N'' / ``Figure N''. It already correctly caught \code{1.73x} inside a Table~9 context in testing and blocked it --- but if the papers reference tables/figures another way (e.g.\ \code{\textbackslash cref}, \code{Tab.~\textbackslash ref\{...\}}), extend those two patterns.
  \item \textbf{Numeric extraction was buggy on first pass and has been fixed and re-tested:} the original regex missed \code{"1.73x"} (a trailing letter broke the word-boundary check) --- it now explicitly handles \code{\%}, \code{x}/\texttimes, and \code{s}/\code{sec} suffixes. Re-verify against real \code{\_patched.tex} files before trusting it broadly, since paper prose has more unit variety than the smoke test covered.
  \item \textbf{False positives are cheap, false negatives are not.} The pipeline is deliberately conservative --- e.g.\ it also flagged the bare \code{9} in ``Table 9'' as a separate (blocked) claim, which is harmless noise, not a wrong auto-fix.
\end{itemize}

\section{Still open}

\begin{itemize}[leftmargin=1.2em]
  \item Real \code{paper\_targets.json} schema (item 6 in Section~5) --- blocks Tier 1's section$\to$results resolution from working correctly until fixed.
  \item Whether \code{audit\_registry.json}'s \code{resolved}/\code{open} items (2, 3, 6, 7, 8, 9, 11) should also get Tier-1 numeric checks run against them now, or be left for manual/Tier-2 handling since several are pure text fixes with no numeric component.
\end{itemize}

\end{document}
